\documentclass[11pt]{amsart}
\usepackage[T1]{fontenc}
\usepackage{lmodern,microtype,amsmath,amssymb,amsthm,mathtools,booktabs}
\usepackage[margin=1.1in]{geometry}
\usepackage[hidelinks]{hyperref}
\usepackage{xurl}

\newtheorem{theorem}{Theorem}[section]
\newtheorem{lemma}[theorem]{Lemma}
\newtheorem{proposition}[theorem]{Proposition}
\newtheorem{corollary}[theorem]{Corollary}
\theoremstyle{remark}
\newtheorem{remark}[theorem]{Remark}
\numberwithin{equation}{section}

\newcommand{\Z}{\mathbb{Z}}
\newcommand{\TV}{\mathrm{TV}}
\newcommand{\tmix}{t_{\mathrm{mix}}}
\newcommand{\Dir}{\mathcal{E}}

\hypersetup{pdftitle={Adding one permutation to a transposition and a long cycle},pdfauthor={Yitzchak Shmalo}}

\title[One added permutation]{Adding one permutation to a transposition and a long cycle}
\author{Yitzchak Shmalo}
\address{Einstein Institute of Mathematics, The Hebrew University of Jerusalem, Jerusalem, Israel}
\email{yitzchak.shmalo@gmail.com}
\date{November 2026}
\subjclass[2020]{60J10, 20B30, 05C25, 60B15}
\keywords{random walk on the symmetric group, mixing time, comparison of Dirichlet forms, Cayley graph, card shuffling}

\begin{document}

\begin{abstract}
For odd $n$, the random walk on the symmetric group $S_n$ generated by the transposition $(1,2)$ and the $n$-cycle $(1,2,\dots,n)$ needs order $n^3\log n$ steps to mix. Diaconis asked whether adding one more permutation to this generating set can speed up the walk appreciably. We show that adding the doubling map $x\mapsto 2x \bmod n$ of $\Z/n\Z$ reduces the mixing time to $O(n\log^3 n)$. Every transposition is then a word of length $O(\log n)$ in the five moves, and a comparison with the random transposition walk that keeps track of the whole spectrum gives the bound. In the other direction, whatever permutation is added, the walk needs order $n\log n$ steps. So the best added permutation brings the mixing time from order $n^3\log n$ to $n^{1+o(1)}$; its exact order, and which permutation attains it, remain open.
\end{abstract}

\maketitle

\section{Introduction}\label{sec:intro}

Label the cards of a deck of odd size $n\ge 5$ by $\Z/n\Z$. Let $\tau=(0,1)$ be the transposition of the first two labels and let $c(x)=x+1$ be the $n$-cycle. The random walk on $S_n$ that at each step multiplies by $\tau$, $c$ or $c^{-1}$ needs order $n^3\log n$ steps to become random. Diaconis and Saloff-Coste proved an upper bound of this order by comparison with the random transposition walk~\cite[Theorem~2]{DSC}, and Wilson proved a lower bound of the same order for the walk with these moves, which he calls the symmetrized Rudvalis shuffle~\cite[Theorem~1]{Wilson}. In a survey of questions on Markov chain Monte Carlo, Diaconis asked which element should be added to a fixed generating set of a group to speed up mixing most, and in particular whether, for $S_n$ with $n$ odd and the generating set $\{\tau,c\}$, some added element speeds up the walk appreciably~\cite[Question~3]{Diaconis13}.

We answer the second question in the affirmative. Let
\[
 \sigma(x)=2x \bmod n,\qquad x\in\Z/n\Z .
\]
Since $n$ is odd, $\sigma$ is a permutation. Put
\[
 T=\{\tau,c,c^{-1},\sigma,\sigma^{-1}\},\qquad h=\lceil\log_2 n\rceil .
\]
For $n\ge5$ these five permutations are distinct and none is the identity: $\sigma$ fixes $0$ while $c^{\pm1}$ do not, $\sigma$ moves $1$ to $2$ while $\tau$ does not, and $\sigma=\sigma^{-1}$ would force $n\mid 3$. The \emph{half-lazy walk} $P_T$ stays put with probability $1/2$ and otherwise multiplies on the right by an element of $T$ chosen uniformly:
\[
 P_T(g,gs)=\tfrac1{10}\quad(s\in T),\qquad P_T(g,g)=\tfrac12 .
\]
It is symmetric, its stationary distribution is the uniform distribution $U$ on $S_n$, and its distance to $U$ does not depend on the starting point. Write $P_T^t(e,\cdot)$ for the law after $t$ steps started at the identity $e$, and
\[
 \tmix(\varepsilon)=\min\{t\ge0:\ \|P_T^t(e,\cdot)-U\|_{\TV}\le\varepsilon\}.
\]

\begin{theorem}\label{thm:main}
Let $n\ge5$ be odd. Every transposition of $S_n$ is a product of at most $12h-9$ elements of $T$, and for every $\varepsilon\in(0,1)$,
\[
 \tmix(\varepsilon)\le\Bigl\lceil 720\,h^2(n-1)\Bigl[2\log n+\log\Bigl(1+\frac1{4\varepsilon^2}\Bigr)\Bigr]\Bigr\rceil .
\]
In particular $\tmix(1/4)=O(n\log^3 n)$. The Cayley graph of $S_n$ with respect to $T$ has diameter of order $n\log n$.
\end{theorem}

The lower bound holds for every added permutation. For $\pi\in S_n$ let $T_\pi$ be the set of distinct non-identity elements of $\{\tau,c,c^{-1},\pi,\pi^{-1}\}$, and let $P_{T_\pi}$ be the half-lazy walk that holds with probability $1/2$ and otherwise moves by an element of $T_\pi$ chosen uniformly.

\begin{proposition}\label{prop:lower}
Let $n\ge3$ and $\pi\in S_n$. The ball of radius $r$ about $e$ in the Cayley graph of $S_n$ with respect to $T_\pi$ has at most $(5\cdot4^r-2)/3$ elements. For every fixed $\varepsilon\in(0,1)$ the mixing time of $P_{T_\pi}$ is at least $(2-o(1))\log(n!)/\log4$, uniformly in $\pi$.
\end{proposition}

That a walk driven by a bounded number of generators needs order $n\log n$ steps is the counting bound recorded by Diaconis and Saloff-Coste~\cite[Remark after Theorem~7]{DSC}. Proposition~\ref{prop:lower} records the constant that comes from $\tau$ being an involution. Together the two results give the following.

\begin{corollary}\label{cor:order}
As $n\to\infty$ through odd integers,
\[
 \frac{2-o(1)}{\log 4}\,n\log n\ \le\ \min_{\pi\in S_n}\tmix^{(\pi)}(1/4)\ \le\ C\,n\log^3 n,
\]
where $\tmix^{(\pi)}$ is the mixing time of $P_{T_\pi}$ and $C$ is an absolute constant. Adding one permutation lowers the exponent of $n$ in the mixing time from $3$ to $1$.
\end{corollary}

\subsection*{Related work}
Diaconis and Saloff-Coste~\cite[Theorem~7]{DSC} proved an $O(n\log^3 n)$ bound for a walk with three generators of a different kind: $n$ even, $S_n$ acting on $\Z/(n-1)\Z\cup\{\infty\}$, the two maps $x\mapsto 2x$ and $x\mapsto 2x+1$ fixing $\infty$, and the transposition $(0,\infty)$. These generators come from the work of Babai, Hetyei, Kantor, Lubotzky and Seress on short words in finite groups~\cite{BHKLS}. Our proof follows the same comparison strategy. What is new is the construction for the generating set that Diaconis prescribed: the walk must keep $\tau$ and $c$, and only one permutation may be added. The relation $\sigma c\sigma^{-1}=c^2$ lets the walk move by powers of two along the cycle, and Lemma~\ref{lem:transposition} turns this into short words for all transpositions.

Kassabov~\cite{Kassabov} constructed bounded generating sets of $S_n$ whose Cayley graphs are expanders, so the corresponding walks mix in $O(\log n!)=O(n\log n)$ steps. Those sets are not of the form $\{\tau,c,\pi\}$, and they do not say how fast the best walk of this form mixes. Babai, Kantor and Lubotzky~\cite{BKL} showed that every finite simple group has a generating set of bounded size whose Cayley graph has logarithmic diameter, and for random pairs of generators of $S_n$ Helfgott, Seress and Zuk~\cite{HSZ} bounded the diameter and the mixing time. Chatterjee and Diaconis~\cite{CD} showed that composing a Markov chain with a deterministic bijection of the state space can speed it up greatly. Their bijections act on $S_n$ as a set and need not be multiplications by a group element, so they do not answer Question~3 of~\cite{Diaconis13}. The comparison method itself is due to Diaconis and Saloff-Coste~\cite{DSC}, and the spectrum of the random transposition walk was found by Diaconis and Shahshahani~\cite{DS81}.

The proofs below are self-contained apart from the classical facts on the representations of $S_n$ used in Section~\ref{sec:mixing}. The rest of the paper is organised as follows.

\begin{center}\small
\begin{tabular}{@{}lll@{}}\toprule
Section & What it proves & What it uses \\ \midrule
\ref{sec:words} & short words for $c^i$ and for all transpositions; the diameter & the relation $\sigma c\sigma^{-1}=c^2$ \\
\ref{sec:mixing} & the comparison inequality and Theorem~\ref{thm:main} & Section~\ref{sec:words}; the spectrum of $Q$~\cite{DS81} \\
\ref{sec:lower} & Proposition~\ref{prop:lower} and Corollary~\ref{cor:order} & counting reduced words \\
\ref{sec:small} & exact mixing times for $n=5,7,9$ & computation \\
\ref{sec:open} & two open questions & \\ \bottomrule
\end{tabular}
\end{center}

\subsection*{Conventions}
Products of permutations act on the right-most factor first: for permutations $g$ and $k$, $(gk)(x)=g(k(x))$. A \emph{word} of length $\ell$ for a permutation $g$ is an expression $g=s_1s_2\cdots s_\ell$ with every $s_j\in T$. We write $\log$ for the natural logarithm. Constants implied by $O(\cdot)$ and $o(\cdot)$ are absolute. A table of notation is in Appendix~\ref{app:notation}.

\section{Short words}\label{sec:words}

Since $\sigma c\sigma^{-1}(\sigma(x))=\sigma(x+1)=2x+2=c^2(\sigma(x))$,
\begin{equation}\label{eq:doubling}
 \sigma c\sigma^{-1}=c^2,\qquad\text{and hence}\qquad \sigma^{j}c\,\sigma^{-j}=c^{2^j}\quad(j\ge0).
\end{equation}

\begin{lemma}\label{lem:powers}
For $0<i<n$ the cycle power $c^i$ has a word of length at most $3h-2$.
\end{lemma}

\begin{proof}
By~\eqref{eq:doubling}, $c^{2q+b}=\sigma c^{q}\sigma^{-1}c^{b}$ for every integer $q\ge1$ and every bit $b\in\{0,1\}$. Starting from the word $c$ for $c^1$ and reading the binary digits of $i$ after the leading one, each further digit adds at most three letters. If $i$ has $k$ binary digits, the word has length at most $1+3(k-1)=3k-2$, and $k\le h$ because $i<n\le 2^h$.
\end{proof}

\begin{lemma}\label{lem:transposition}
For $i\in\Z/n\Z$ and $1\le d<n$ the transposition $(i,i+d)$ has a word of length at most $12h-9$.
\end{lemma}

\begin{proof}
Write $d=2^{j_1}+\dots+2^{j_r}$ with $0\le j_1<\dots<j_r\le h-1$, and let $a_\ell=\sum_{u<\ell}2^{j_u}$ for $1\le\ell\le r+1$, so that $a_1=0$ and $a_{r+1}=d$. The numbers $a_1<a_2<\dots<a_{r+1}$ lie in $[0,d]$ and $d<n$, so they are distinct residues. Put $A_\ell=c^{a_\ell}\sigma^{j_\ell}$. Since $A_\ell(0)=a_\ell$ and $A_\ell(1)=a_\ell+2^{j_\ell}=a_{\ell+1}$,
\[
 e_\ell:=A_\ell\tau A_\ell^{-1}=(a_\ell,a_{\ell+1}).
\]
The transpositions $e_1,\dots,e_r$ are the edges of the path $a_1,a_2,\dots,a_{r+1}$, so
\[
 (0,d)=e_1e_2\cdots e_{r-1}\,e_r\,e_{r-1}\cdots e_1=B\tau B^{-1},\qquad B=e_1e_2\cdots e_{r-1}A_r .
\]
Expanding $e_\ell=A_\ell\tau A_\ell^{-1}$ gives $B=A_1\tau(A_1^{-1}A_2)\tau(A_2^{-1}A_3)\cdots\tau(A_{r-1}^{-1}A_r)$, and by~\eqref{eq:doubling}
\[
 A_\ell^{-1}A_{\ell+1}=\sigma^{-j_\ell}c^{2^{j_\ell}}\sigma^{j_{\ell+1}}=c\,\sigma^{j_{\ell+1}-j_\ell}.
\]
Hence
\[
 B=\sigma^{j_1}\,\tau c\,\sigma^{j_2-j_1}\,\tau c\,\sigma^{j_3-j_2}\cdots\tau c\,\sigma^{j_r-j_{r-1}},
\]
a word of length $j_r+2(r-1)\le 3h-3$, and $(0,d)=B\tau B^{-1}$ has a word of length at most $6h-5$. Finally $(i,i+d)=c^{i}(0,d)c^{-i}$, and Lemma~\ref{lem:powers} gives a word of length at most $(6h-5)+2(3h-2)=12h-9$.
\end{proof}

Every permutation of $n$ letters is a product of at most $n-1$ transpositions, so Lemma~\ref{lem:transposition} bounds the diameter of the Cayley graph by $(12h-9)(n-1)=O(n\log n)$. The lower bound of the same order follows from the ball estimate in Proposition~\ref{prop:lower}, since a ball containing all of $S_n$ has $n!$ elements. This proves the last assertion of Theorem~\ref{thm:main}.

\section{The mixing bound}\label{sec:mixing}

Let $N=n!$ and $M=\binom n2$. For a function $f$ on $S_n$ and a permutation $s$ put
\[
 D_s(f)=\frac1N\sum_{g\in S_n}\bigl(f(gs)-f(g)\bigr)^2 .
\]
Let $Q$ be the half-lazy random transposition walk, which holds with probability $1/2$ and otherwise multiplies by one of the $M$ transpositions chosen uniformly. The Dirichlet forms of $P_T$ and $Q$ are
\[
 \Dir_{P_T}(f)=\frac1{20}\sum_{s\in T}D_s(f),\qquad \Dir_{Q}(f)=\frac1{4M}\sum_{u}D_u(f),
\]
where $u$ runs over the transpositions.

\begin{lemma}\label{lem:comparison}
Let $A=720h^2$. Then $\Dir_Q(f)\le A\,\Dir_{P_T}(f)$ for every $f$.
\end{lemma}

\begin{proof}
Fix a transposition $u$ and the word $u=s_1\cdots s_\ell$ of Lemma~\ref{lem:transposition}, so $\ell\le L:=12h$. Writing $g_k=gs_1\cdots s_k$, we have $f(gu)-f(g)=\sum_{k=1}^{\ell}\bigl(f(g_k)-f(g_{k-1})\bigr)$, and the Cauchy--Schwarz inequality gives $(f(gu)-f(g))^2\le\ell\sum_k(f(g_{k-1}s_k)-f(g_{k-1}))^2$. Averaging over $g$, and using that $g\mapsto g_{k-1}$ is a bijection of $S_n$,
\[
 D_u(f)\le\ell\sum_{k=1}^{\ell}D_{s_k}(f)\le L^2\sum_{s\in T}D_s(f),
\]
because each $s\in T$ occurs at most $\ell\le L$ times in the word. Summing over the $M$ transpositions,
\[
 \Dir_Q(f)\le\frac{L^2}{4}\sum_{s\in T}D_s(f)=5L^2\,\Dir_{P_T}(f)=720h^2\,\Dir_{P_T}(f).\qedhere
\]
\end{proof}

Both $P_T$ and $Q$ are symmetric, and since they hold with probability $1/2$ their eigenvalues lie in $[0,1]$. Let $0=\beta_0<\beta_1\le\dots\le\beta_{N-1}$ be the eigenvalues of $I-P_T$ and $0=\mu_0<\mu_1\le\dots\le\mu_{N-1}$ those of $I-Q$, with multiplicity. By Lemma~\ref{lem:comparison} and the min--max characterisation of eigenvalues, $\beta_k\ge\mu_k/A$ for every $k$.

The eigenvalues of $Q$ are classical~\cite{DS81}. The irreducible representations of $S_n$ are indexed by the partitions $\lambda$ of $n$; let $d_\lambda$ be the dimension and $\lambda'$ the conjugate partition. The average of the transpositions acts on the $\lambda$-isotypic part of the regular representation, which has dimension $d_\lambda^2$, as multiplication by
\[
 r_\lambda=\frac{\sum_i\lambda_i^2-\sum_i(\lambda'_i)^2}{n(n-1)},
\]
so $I-Q$ has the eigenvalue $(1-r_\lambda)/2$ with multiplicity $d_\lambda^2$. Write $j=n-\lambda_1$. Since $\sum_i\lambda_i^2\le\lambda_1\sum_i\lambda_i=n\lambda_1$ and $\sum_i(\lambda'_i)^2\ge\sum_i\lambda'_i=n$,
\begin{equation}\label{eq:mu}
 \frac{1-r_\lambda}{2}\ \ge\ \frac{n(n-1)-n\lambda_1+n}{2n(n-1)}=\frac{j}{2(n-1)}.
\end{equation}
For the multiplicities: a standard Young tableau of shape $\lambda$ is determined by the set of $j$ labels outside its first row together with the standard tableau, of shape $\bar\lambda=(\lambda_2,\lambda_3,\dots)$, that they form after relabelling in increasing order. Hence $d_\lambda\le\binom nj d_{\bar\lambda}$, and since $\sum_{\nu\vdash j}d_\nu^2=j!$,
\begin{equation}\label{eq:mult}
 \sum_{\lambda:\ \lambda_1=n-j}d_\lambda^2\ \le\ \binom nj^2 j!\ \le\ n^{2j}.
\end{equation}

\begin{proof}[Proof of Theorem~\ref{thm:main}]
The word bound and the diameter were proved in Section~\ref{sec:words}. For the mixing bound, the walk is a symmetric random walk on a group, so $P_T^{2t}(g,g)=P_T^{2t}(e,e)$ for all $g$ and
\[
 \chi^2\bigl(P_T^t(e,\cdot),U\bigr):=N\sum_g P_T^t(e,g)^2-1=N\,P_T^{2t}(e,e)-1=\operatorname{tr}(P_T^{2t})-1=\sum_{k\ge1}(1-\beta_k)^{2t}.
\]
Using $0\le1-\beta_k\le e^{-\beta_k}$, then $\beta_k\ge\mu_k/A$, and then~\eqref{eq:mu} and~\eqref{eq:mult},
\[
 \chi^2\bigl(P_T^t(e,\cdot),U\bigr)\le\sum_{k\ge1}e^{-2t\mu_k/A}\le\sum_{j=1}^{n-1}n^{2j}e^{-tj/(A(n-1))}=\sum_{j=1}^{n-1}x^j,\qquad x=n^2e^{-t/(A(n-1))}.
\]
If $t\ge A(n-1)\bigl[2\log n+\log(1+1/(4\varepsilon^2))\bigr]$ then $x\le 4\varepsilon^2/(1+4\varepsilon^2)$, so $\sum_{j\ge1}x^j\le x/(1-x)\le 4\varepsilon^2$. By the Cauchy--Schwarz inequality,
\[
 \|P_T^t(e,\cdot)-U\|_{\TV}=\frac12\sum_g\bigl|P_T^t(e,g)-\tfrac1N\bigr|\le\frac12\sqrt{\chi^2\bigl(P_T^t(e,\cdot),U\bigr)}\le\varepsilon.\qedhere
\]
\end{proof}

\begin{remark}
Keeping only the spectral gap $\beta_1\ge 1/(A(n-1))$ and using the minimum stationary probability $1/n!$ would give a bound of order $n^2\log^3 n$. The proof above avoids this loss because it compares the whole spectrum and controls the multiplicities through~\eqref{eq:mult}. No attempt has been made to optimise the constant $720$.
\end{remark}

\section{The lower bound}\label{sec:lower}

\begin{proof}[Proof of Proposition~\ref{prop:lower}]
The set $T_\pi$ is closed under inverses and has at most five elements, and $\tau=\tau^{-1}$. An element at distance $r\ge1$ from $e$ has a word of length $r$ in which no letter is followed by its inverse. There are at most five choices for the first letter and at most four for each later one, so the ball $B(r)$ of radius $r$ has at most $1+5\sum_{k=0}^{r-1}4^k=(5\cdot4^r-2)/3$ elements.

Let $R=\log_4(n!)$, fix $\delta\in(0,1)$, and let $t=\lfloor(2-\delta)R\rfloor$ and $r=\lfloor(1-\delta/4)R\rfloor$. After $t$ steps the walk is at a product of $K_t$ elements of $T_\pi$, where $K_t$ is the number of steps at which it did not hold, a binomial variable with parameters $t$ and $1/2$. Its mean is at most $(1-\delta/2)R$ and its variance at most $t/4\le R/2$, so by Chebyshev's inequality
\[
 \mathbb P(K_t>r)\le\frac{R/2}{(\delta R/4-1)^2}\to0
\]
as $n\to\infty$, uniformly in $\pi$. Hence $P_{T_\pi}^t(e,B(r))\ge1-o(1)$, while $U(B(r))\le(5\cdot4^r-2)/(3\,n!)\le\tfrac53\,4^{-\delta R/4}\to0$. The total variation distance at time $t$ therefore tends to $1$, and for any fixed $\varepsilon<1$ the mixing time exceeds $(2-\delta)R-1$ for large $n$. Since $\delta$ was arbitrary, this proves the proposition.
\end{proof}

\begin{proof}[Proof of Corollary~\ref{cor:order}]
The lower bound is Proposition~\ref{prop:lower}. The upper bound is Theorem~\ref{thm:main} with $\pi=\sigma$, for which $T_\sigma=T$. By Stirling's formula $\log(n!)=(1+o(1))n\log n$, while the base walk on $\{\tau,c,c^{-1}\}$ needs order $n^3\log n$ steps~\cite{DSC,Wilson}.
\end{proof}

\section{Small decks}\label{sec:small}

For $n=5,7,9$ we computed the exact distribution of the half-lazy walk started at $e$, both for the base walk that moves by $\tau$, $c$ or $c^{-1}$ with probability $1/6$ each and for $P_T$, and recorded the first time the total variation distance to $U$ falls to $1/4$ and to $1/100$.

\begin{center}
\begin{tabular}{@{}rrrrrr@{}}\toprule
$n$ & $h$ & base, $\varepsilon=1/4$ & $P_T$, $\varepsilon=1/4$ & base, $\varepsilon=1/100$ & $P_T$, $\varepsilon=1/100$ \\ \midrule
5 & 3 & 29 & 13 & 92 & 44 \\
7 & 3 & 103 & 42 & 281 & 113 \\
9 & 4 & 254 & 74 & 633 & 171 \\ \bottomrule
\end{tabular}
\end{center}

\noindent The ratio of the base mixing time to that of $P_T$ grows with $n$, as Theorem~\ref{thm:main} predicts, although at these sizes it is still small. For $n=5$ we also tried every one of the $120$ permutations $\pi$ as the added generator, with the half-lazy walk $P_{T_\pi}$ defined as in Proposition~\ref{prop:lower}. The smallest time to reach distance $1/4$ is $9$; it is attained by $14$ permutations, one of which is $(\pi(0),\dots,\pi(4))=(1,2,4,0,3)$. The doubling map, with time $13$, is not among them. So doubling is not the best added permutation for small decks, and we do not know which permutation is best for large $n$.

The words of Lemmas~\ref{lem:powers} and~\ref{lem:transposition} were also checked by computer: for every odd $n$ from $5$ to $301$ the words evaluate to the claimed permutations and have the claimed lengths. The longest transposition word has $17$ letters for $n=5$ (bound $27$) and $69$ for $n=101$ (bound $75$). The code is in the repository listed under Data availability.

\section{Questions}\label{sec:open}

Corollary~\ref{cor:order} leaves a factor $\log^2 n$ between the bounds. The same gap appears in Theorem~7 of~\cite{DSC}, where Diaconis and Saloff-Coste asked whether any part of it is necessary. Two questions remain for the walk of Diaconis's Question~3.

\begin{enumerate}
\item The order of $\min_\pi\tmix^{(\pi)}(1/4)$ as $n\to\infty$ through odd integers.
\item An added permutation $\pi$ that attains this order.
\end{enumerate}

\noindent A permutation $\pi$ for which $\{\tau,c,\pi\}$ generates an expander family would give order $n\log n$; the computation for $n=5$ suggests that the best choice is not the doubling map.

\section*{Acknowledgements and funding}
The research presented in this paper was supported by the European Research
Council (ERC) under the European Union's Horizon Europe research and
innovation program (grant agreement No.~101041711), by the Simons
Foundation as part of the Collaboration on the Mathematical and Scientific
Foundations of Deep Learning, by Heights Labs, by Convex Nexus Capital, by
the Israel Science Foundation (grant number 2258/19), and by the Israel
Science Foundation (ISF Grant 4101/25). Computational resources were
provided in part by Google Cloud through its education credits program
(credit 530550103).
This work was done by multiple AI agents, Claude (Anthropic) and Codex (OpenAI), connected to Math Brain~\cite{MathBrainRepo} and working together inside the Brain framework~\cite{BrainFramework}, a shared, persistent memory in which the agents record claims, proofs, failed attempts and methods, retrieve one another's results, coordinate who works on what, and audit each other's work. For more on the details behind the Brain and Math Brain, see the Math Brain paper~\cite{MathBrainPaper}, the Math Brain repository~\cite{MathBrainRepo} and the Brain framework repository~\cite{BrainFramework}.

\section*{Declaration of competing interest}
The author holds equity in Heights Labs and in Convex Nexus Capital and
has a management role at Convex Nexus Capital. These entities are included
in the funding declaration above.

\section*{Declaration of generative AI and AI-assisted technologies}
The paper comes out of Math Brain, an ongoing project on the use of AI systems for research mathematics. In November 2026 Codex (OpenAI, GPT-6.1) found the proofs of Theorem~\ref{thm:main} and Proposition~\ref{prop:lower}, a second Codex instance reviewed the proofs, and Claude Opus~5.5 (Anthropic), working in Claude Code, wrote the paper and the verification code. The author directed the work and is responsible for the content of the paper and for decisions concerning its dissemination.

\section*{Data availability}
No data were generated or analysed. The \LaTeX{} source of the paper and the code that checks the words of Section~\ref{sec:words} and computes the table of Section~\ref{sec:small} are available at \url{https://github.com/yspennstate/one-generator-shuffle}.

\appendix
\section{Notation}\label{app:notation}
Symbols are listed in the order in which they first appear.

\begin{center}\small
\begin{tabular}{@{}lp{0.62\textwidth}l@{}}\toprule
Symbol & Meaning & First appears \\ \midrule
$n$ & the number of cards, odd in Theorem~\ref{thm:main} & Section~\ref{sec:intro} \\
$\tau$ & the transposition $(0,1)$ & Section~\ref{sec:intro} \\
$c$ & the $n$-cycle $x\mapsto x+1$ on $\Z/n\Z$ & Section~\ref{sec:intro} \\
$\sigma$ & the doubling map $x\mapsto 2x\bmod n$ & Section~\ref{sec:intro} \\
$T$ & the move set $\{\tau,c,c^{-1},\sigma,\sigma^{-1}\}$ & Section~\ref{sec:intro} \\
$h$ & $\lceil\log_2 n\rceil$ & Section~\ref{sec:intro} \\
$P_T$ & the half-lazy walk with moves in $T$ & Section~\ref{sec:intro} \\
$e$, $U$ & the identity permutation and the uniform distribution on $S_n$ & Section~\ref{sec:intro} \\
$\tmix(\varepsilon)$ & the first time the walk is within $\varepsilon$ of $U$ in total variation & Section~\ref{sec:intro} \\
$\pi$, $T_\pi$, $P_{T_\pi}$ & an added permutation, the distinct non-identity elements of $\{\tau,c,c^{-1},\pi,\pi^{-1}\}$, and the half-lazy walk with these moves & Section~\ref{sec:intro} \\
$\tmix^{(\pi)}$ & the mixing time of $P_{T_\pi}$ & Corollary~\ref{cor:order} \\
$a_\ell$, $j_\ell$, $A_\ell$, $e_\ell$, $B$ & the partial sums of $d$, its binary exponents, and the permutations built from them & Lemma~\ref{lem:transposition} \\
$N$, $M$ & $n!$ and $\binom n2$ & Section~\ref{sec:mixing} \\
$D_s(f)$ & the mean square change of $f$ under right multiplication by $s$ & Section~\ref{sec:mixing} \\
$Q$, $u$ & the half-lazy random transposition walk, and a transposition & Section~\ref{sec:mixing} \\
$\Dir_{P_T}$, $\Dir_Q$ & the Dirichlet forms of $P_T$ and $Q$ & Section~\ref{sec:mixing} \\
$A$, $L$ & the comparison constant $720h^2$ and the word bound $12h$ & Lemma~\ref{lem:comparison} \\
$\beta_k$, $\mu_k$ & the eigenvalues of $I-P_T$ and of $I-Q$ & Section~\ref{sec:mixing} \\
$\lambda$, $\lambda'$, $\bar\lambda$, $d_\lambda$ & a partition of $n$, its conjugate, $\lambda$ without its first row, and the dimension of the representation & Section~\ref{sec:mixing} \\
$r_\lambda$ & the character ratio of a transposition at $\lambda$ & Section~\ref{sec:mixing} \\
$\chi^2$ & the chi-square distance to $U$ & Section~\ref{sec:mixing} \\
$B(r)$, $R$, $K_t$ & the ball of radius $r$, $\log_4(n!)$, and the number of non-holding steps & Section~\ref{sec:lower} \\
\bottomrule
\end{tabular}
\end{center}

\begin{thebibliography}{99}
\bibitem{BKL} L.~Babai, W.~M.~Kantor and A.~Lubotzky, \emph{Small-diameter Cayley graphs for finite simple groups}, European J.\ Combin.\ \textbf{10} (1989), 507--522. doi:10.1016/S0195-6698(89)80067-8.
\bibitem{BHKLS} L.~Babai, G.~Hetyei, W.~M.~Kantor, A.~Lubotzky and \'A.~Seress, \emph{On the diameter of finite groups}, in: Proceedings of the 31st Annual Symposium on Foundations of Computer Science, IEEE, 1990, 857--865. doi:10.1109/FSCS.1990.89608.
\bibitem{CD} S.~Chatterjee and P.~Diaconis, \emph{Speeding up Markov chains with deterministic jumps}, Probab.\ Theory Related Fields \textbf{178} (2020), 1193--1214. doi:10.1007/s00440-020-01006-4.
\bibitem{Diaconis13} P.~Diaconis, \emph{Some things we've learned (about Markov chain Monte Carlo)}, Bernoulli \textbf{19} (2013), 1294--1305. doi:10.3150/12-BEJSP09.
\bibitem{DSC} P.~Diaconis and L.~Saloff-Coste, \emph{Comparison techniques for random walk on finite groups}, Ann.\ Probab.\ \textbf{21} (1993), 2131--2156. doi:10.1214/aop/1176989013.
\bibitem{DS81} P.~Diaconis and M.~Shahshahani, \emph{Generating a random permutation with random transpositions}, Z.\ Wahrsch.\ Verw.\ Gebiete \textbf{57} (1981), 159--179. doi:10.1007/BF00535487.
\bibitem{HSZ} H.~A.~Helfgott, \'A.~Seress and A.~Zuk, \emph{Random generators of the symmetric group: diameter, mixing time and spectral gap}, J.\ Algebra \textbf{421} (2015), 349--368. doi:10.1016/j.jalgebra.2014.08.033.
\bibitem{Kassabov} M.~Kassabov, \emph{Symmetric groups and expander graphs}, Invent.\ Math.\ \textbf{170} (2007), 327--354. doi:10.1007/s00222-007-0065-y.
\bibitem{Wilson} D.~B.~Wilson, \emph{Mixing time of the Rudvalis shuffle}, Electron.\ Commun.\ Probab.\ \textbf{8} (2003), 77--85. doi:10.1214/ECP.v8-1071.
\bibitem{MathBrainRepo} Y.~Shmalo, \emph{Math Brain: public snapshot, code and tests}, 2026. doi:10.5281/zenodo.24150837.
\bibitem{MathBrainPaper} Y.~Shmalo, \emph{Math Brain: shared, auditable memory for multi-agent mathematical research}, October 2026, \url{https://github.com/yspennstate/math-brain-public/tree/9eb386da3dc67da2fbb4a78a9d97baaaa1f79b9a/papers/math-brain-2026-10-24}.
\bibitem{BrainFramework} Y.~Shmalo, \emph{Brain Framework}, public software repository, 2026, commit 5629e64, \url{https://github.com/yspennstate/brain-framework-public/tree/5629e644cdd510173ae97e9aa92642dabd22dc57}.
\end{thebibliography}
\end{document}
